\section{The Encoding Is Abstract, Steerable, and Portable}
\label{sec:steer}

\Cref{sec:what,sec:heads} showed the goal is decodable from a sparse set of heads. We now show it is a
genuinely \emph{abstract} representation of intent: transplanting it \emph{steers} a fresh context to
pursue the intended sub-goal, and it does so across programming languages, across model families, and
even outside code. Throughout, a \emph{source} run $\M(p_\src)$ supplies the \goalheads{} activations we
transplant into a \emph{destination} run $\M(p_\dst)$; we report the oracle-judged \textbf{causal
trigger rate} (does the destination's next computation switch to the intended sub-goal
$g_\mathrm{targ}$?) and, where applicable, the \textbf{readout-shift} (does \goalscope on the patched
destination now verbalize $g_\mathrm{targ}$?) and \textbf{gaze-shift} (do the \goalheads now look back
to $g_\mathrm{targ}$?).

\subsection{Transplanting the encoding steers the next sub-goal}
\label{sec:steer:transplant}

We construct source/destination pairs that share a task but are about to pursue \emph{different} next
sub-goals ($g_\src \neq g_\dst$). Transplanting the source \goalheads{} states at the pivot makes the
destination pursue $g_\src$ instead of its native $g_\dst$
({\color{red}\Cref{fig:steer-transplant}}). Crucially this is not a generic ``perturb-and-derail'': the
destination continues to write correct, coherent code, but for the transplanted sub-goal. The
readout-shift and gaze-shift track the causal trigger, confirming the intervention moves the intended
latent, not a surface artifact.
\arnab{settle final transplant variant(s) reported here; multi-predicate ordering is the cleanest
``which-goal'' contrast.}

\begin{figure}[t]
\centering
\begin{subfigure}{0.48\linewidth}
\includegraphics[width=\linewidth]{Figures/axis_q_multipred.png}
\caption{query/OV transplant}
\end{subfigure}\hfill
\begin{subfigure}{0.48\linewidth}
\includegraphics[width=\linewidth]{Figures/axis_kv_multipred.png}
\caption{key/value transplant}
\end{subfigure}
\caption{Transplanting the goal encoding steers the destination to the intended sub-goal (placeholders).}
\label{fig:steer-transplant}
\end{figure}

\subsection{Portable across programming languages}
\label{sec:steer:crosslang}

The transplant works \emph{across the language boundary}: a goal cached from a Python context triggers
the same sub-goal when injected into a C++ context, and vice versa. Because the source's target
implementation is never present in the destination's language, cross-language transfer is direct evidence
that the transplanted representation encodes the abstract sub-task, not its surface form. We report
per-direction rates (py$\to$cpp, cpp$\to$py) against a same-language upper bound in \Cref{app:cpp}.

\subsection{Portable across model families and training regimes}
\label{sec:steer:crossmodel}

The phenomenon is not idiosyncratic to one model. We replicate the localization and \goalscope read-out
across (a) \emph{scale}, (b) \emph{code fine-tuning}, and (c) \emph{RL post-training}, using matched
head budgets per model (\Cref{fig:cross-model}; OLMo details in \Cref{app:olmo}). A compact set of goal
heads carrying a decodable, steerable next-sub-goal encoding appears across all of them.

\begin{figure}[t]
\centering
\begin{subfigure}{0.32\linewidth}
\includegraphics[width=\linewidth]{Figures/goalscope_axis_c_scale.png}
\caption{scale}
\end{subfigure}\hfill
\begin{subfigure}{0.32\linewidth}
\includegraphics[width=\linewidth]{Figures/goalscope_axis_c_code_ft.png}
\caption{code fine-tuning}
\end{subfigure}\hfill
\begin{subfigure}{0.32\linewidth}
\includegraphics[width=\linewidth]{Figures/goalscope_axis_c_rl.png}
\caption{RL post-training}
\end{subfigure}
\caption{The decodable/steerable goal encoding generalizes across model families and training regimes
(placeholders).}
\label{fig:cross-model}
\end{figure}

\subsection{Beyond code: few-shot and sequential planning}
\label{sec:steer:noncode}

Finally, the account is not specific to code. In few-shot in-context-learning prompts, \goalscope
verbalizes the task the model is about to apply; and in non-code sequential-planning settings (e.g.\
following a recipe or an itinerary) we find qualitative read-out and steering signals of the same kind.
We report these out-of-domain results in \Cref{app:non-code}.
\arnab{quantify the ICL read-out (we said this is likely doable); keep planning as qualitative for now.}
